\section{Before Jerome: The Old Latin Bible}

Latin Christianity did not wait for Jerome. By the time he took up his pen,
Latin Scripture was some two centuries old, and its state was the problem
he was hired to solve.

The witnesses agree on the disorder. Augustine, writing his handbook of
Christian learning around 396, says that while the translators of the
Scriptures from Hebrew into Greek can be counted, ``the Latin translators
are out of all number,'' for ``in the early days of the faith every man who
happened to get his hands upon a Greek manuscript, and who thought he had
any knowledge, were it ever so little, of the two languages, ventured upon
the work of translation.''\endnote{Augustine, \work{On Christian Doctrine}
2.11.16, in the J.~F. Shaw translation (NPNF series~1, vol.~2), checked at
New Advent. Modern scholarship cautions that the remark concerns the Old
Testament in context and has ``long been overapplied'' (Houghton,
\work{Latin New Testament}, ch.~1).} Jerome, in the preface that opens his
own revision, is blunter: if the faith is to rest on the Latin version,
let his critics answer which one---``for they are nearly as many as the
books!''\endnote{Jerome, preface to the Gospels (\textit{Novum opus}),
addressed to Damasus, in the public-domain translation of Kevin~P.
Edgecomb at tertullian.org (``If indeed faith is administered by the
Latin version, they might respond by which, for they are nearly as many
as the books!''). The Latin tag is \textit{tot sunt paene quot
codices}.}

Two corrections keep this famous chaos in proportion. First, the disorder
the Fathers describe is largely a disorder of copies and revisions, not
necessarily of independent origins: the editors of the modern Beuron
\work{Vetus Latina} edition have concluded, book by book, that a single
Latin translation appears to underlie the surviving Old Latin evidence for
each biblical book, subsequently revised, corrected toward Greek copies,
and cross-contaminated almost without
restraint.\endnote{Houghton, \work{Latin New Testament}, ch.~1, reporting
the consensus of the \work{Vetus Latina} editors (Fischer, Thiele, Burton,
and others cited there) for the New Testament and, with Haelewyck, for the
Old. The shared renderings of rare Greek words across all witnesses are
among the arguments. The conclusion is a modern reconstruction (class~K),
book-specific, and does not exclude lost early parallel versions.}
``Old Latin'' or \work{Vetus Latina} is thus a modern collective label for
everything earlier than or independent of the Vulgate revisions---a
plurality of text forms, not one ancient edition. Second, the base text was
Greek throughout: for the Old Testament the Septuagint, not the Hebrew. The
Latin Bible before Jerome was a translation of a translation, and its
Psalter, its Jonah, its Daniel were the Septuagint's.

Augustine also supplies the most disputed sentence in this history. Among
the translations themselves, he writes, ``the Italian (\textit{Itala}) is to be
preferred to the others, for it keeps closer to the words without prejudice
to clearness of expression.''\endnote{Augustine, \work{On Christian
Doctrine} 2.15.22 (Shaw, NPNF series~1, vol.~2), checked at New Advent.}
What the \textit{Itala} was has been argued for three centuries: an
otherwise unattested Italian text-type; a corruption in the manuscripts; or
even, on one influential nineteenth-century theory (Burkitt's), Jerome's
own Vulgate. Current scholarship mostly reads it as a geographical label
for the Italian form of the Old Latin that Augustine had met in
Milan---but the term never recurs, and this account uses it for
Augustine's remark and nothing else.\endnote{The debate is summarized,
with Burkitt 1896 and Quentin 1927 as the outlying positions and the
geographical reading as the present preference, in Houghton, \work{Latin
New Testament}, ch.~1 (class~K). Older handbooks used ``Itala'' for the
whole Old Latin tradition; the usage has been abandoned.}

One term of art completes the setting. When fourth-century writers said
\textit{editio vulgata}, ``the common edition,'' they meant this
Greek-based common text---Jerome himself uses ``common edition'' for the
Septuagint's. The name that now means Jerome's Bible belonged, in his
lifetime, to the Bible he was trying to replace.\endnote{Houghton,
\work{Latin New Testament}, ch.~2; Jerome's usage appears, for example, in
the preface to the books of Solomon (``called in the common edition
Proverbs,'' tr.\ Edgecomb). See the terminal appendix on terminology.}
